% STATUS: QUARANTINED old-method narrative; see ../METHOD_STATUS.md.
\section{引言}
\label{sec:introduction}

小麦器官的数量、空间分布、颜色变化和衰老进程与冠层结构、生长状态及产量形成密切相关 \citep{li2021canopy,li2021lai,dandrifosse2022ear}。高空间分辨率 RGB 成像为这些性状提供了低成本、非破坏性的田间测量手段，但可靠测量仍依赖对麦穗、叶片、茎秆和背景的准确区分。与以麦穗定位为目标的检测任务不同，器官级语义分割必须同时处理叶片交叠、茎秆遮挡、组织衰老、杂草干扰和土壤残体等复杂因素。GWFSS 数据集覆盖多个国家、研究机构、生育期和成像平台，为研究真实田间条件下的小麦器官分割提供了统一基础 \citep{wang2025gwfss}。

困难并不只来自训练样本数量。少标注协议仅提供 99 张像素级标注图像，而无标注图像超过六万张，教师模型生成的伪标签因而承担了大量监督作用。细茎只占少量像素，其边界置信度又显著低于叶片和背景；固定阈值往往把整个不确定掩码删除，使中心位置正确、宽度尚不确定的茎秆监督也随之消失。采集条件的变化进一步放大了这种偏差。不同机构的图像在光照、基因型、生育期、相机视角和地面采样距离上存在明显差异，地域划分下显著低于随机划分的结果表明，未知采集域泛化是实际部署中的核心障碍 \citep{wang2025gwfss}。分辨率则构成另一层制约：放大输入确实有助于识别细茎，但密集覆盖多个尺度会对同一区域反复计算 \citep{cao2025detailminority}，难以满足高通量表型分析的效率要求。

现有研究分别从半监督学习、结构保持和跨域表示等角度缓解上述问题。Mean Teacher 及其后续方法使用指数滑动平均教师为学生提供稳定监督 \citep{tarvainen2017mean,berrada2024guided}；UniMatch、U2PL 和 ST++ 分别从强弱扰动、不可靠像素利用和样本稳定性角度改进伪标签学习 \citep{yang2023unimatch,wang2022u2pl,yang2022stpp}。然而，这些方法主要依据像素、区域或图像层面的置信度处理噪声，没有针对细长器官区分结构位置与边界宽度的不确定性。soft-clDice 等方法能够增强管状目标的拓扑连续性 \citep{shit2021cldice}，但直接将拓扑损失作用于噪声伪标签可能强化错误连接。域泛化方法通过类别记忆或域间对齐学习稳定表示 \citep{kim2022pinmemory}，由全部预测像素构建的原型仍会受到细茎边界噪声污染。

稀缺标注条件下的跨域细茎分割不宜被拆成彼此独立的伪标签筛选、域对齐和高分辨率推理问题。伪标签中稳定的拓扑核心能够为跨域原型提供较干净的语义依据，同一结构表示也能够指示哪些局部区域真正需要更高分辨率。沿着这条信息路径，本文构建了 \METHOD。基础分割器由 BEiTv2-Large、ViT-Adapter 和 Mask2Former 组成 \citep{peng2022beitv2,chen2023vitadapter,cheng2022mask2former}，其教师--学生训练过程同时建立拓扑可靠伪监督和跨域器官原型，推理过程则利用预测中的结构失效信号分配局部计算。

\METHOD 的主要贡献不是简单叠加三个独立部件，而是让拓扑可靠性在训练和推理之间发挥一致作用。\TRPL 通过多尺度与扰动一致性估计，将茎秆监督表示为稳定中心线、可靠区域和不确定边界，避免固定阈值整体丢弃细结构。\TCPM 接收这一稳定核心，以课程日程依次建立有标注和高质量伪标注类别--域记忆，并用逐样本留一域锚点校正可靠区域中偏移最明显的困难特征，使干净的原型来源与有效的梯度位置彼此分离。\BAZR 进一步把结构不确定性转化为推理预算，根据预测熵、骨架端点和内部尺度分歧选择少量高分辨率窗口。少标注跨域评测、完整数据地域划分、针对性消融和精度--效率分析共同验证了这一设计链路。
